%%%%%%%%%%%%%%%%%%%%%%%%%%%%%%%%%%%%%%%%%%%%%%%%%%%%%%%%%%%%%%%%%
%% TargetTrackHisTag manuscript - Bioinformatics Advances
%% Uses the OUP bioinfo.cls document class (download from OUP author
%% guidelines or the Overleaf "Bioinformatics" template gallery entry,
%% then upload bioinfo.cls + natbib.sty alongside this .tex file)
%%%%%%%%%%%%%%%%%%%%%%%%%%%%%%%%%%%%%%%%%%%%%%%%%%%%%%%%%%%%%%%%%%%

\documentclass[unnumsec,webpdf,contemporary,large]{oup-authoring-template}

\journaltitle{Bioinformatics Advances}
\pubyear{2026}
\copyrightyear{2026}
\usepackage{graphicx}
\usepackage{booktabs}
\usepackage{amsmath}

\begin{document}

\title[TargetTrackHisTag]{TargetTrackHisTag: an experimental benchmark reveals a significant but modest sequence-encoded signal for His-tag placement preference}
\author[]{Muhammad Uzair Ashraf\,$^{1}$, Muazzam\,$^{1}$ and Rizwan Hasan Khan\,$^{1,}$\,$^{*}$}
\address{$^{1}$Interdisciplinary Biotechnology Unit, Aligarh Muslim University, Aligarh, 202002, Uttar Pradesh, India}

\corresp{$^\ast$To whom correspondence should be addressed. E-mail: rhkhan.cb@amu.ac.in}

\history{Received on XXXXX; revised on XXXXX; accepted on XXXXX}

\editor{Associate Editor: XXXXXXX}

\abstract{
\textbf{Motivation:} Optimal placement of polyhistidine (His) affinity tags at the N- or C-terminus of recombinant proteins remains empirically determined due to the absence of validated computational predictors and experimentally anchored benchmarks against which sequence-based approaches can be rigorously evaluated.\\
\textbf{Results:} We present TargetTrackHisTag, the first machine-readable benchmark of experimentally derived bilateral His-tag placement outcomes, curated from the Northeast Structural Genomics Consortium (NESG) TargetTrack archive. Parsing 97,418 His-tag construct entries from a 5.4\,GB XML archive, applying alignment-based validation, and filtering at 85\% sequence coverage yielded 156 high-confidence proteins with matched N- and C-terminal constructs and experimentally recorded outcomes. Automated quality audits confirmed 99.4\% genuine His-tag patterns (155/156) and 100\% medium-to-high confidence labels (156/156). Evaluating full-sequence ESM-2 (650M) embeddings with logistic regression against this benchmark using leave-one-out cross-validation across all 156 proteins yielded a statistically significant signal (MCC\,=\,0.168, balanced accuracy\,=\,0.583, permutation $p$\,=\,0.0498, $n$\,=\,156), in contrast to ESMFold pLDDT-derived pseudo-labels (MCC\,=\,$-$0.243), terminal ESM-2 embeddings with XGBoost on pseudo-labelled data evaluated against experimental outcomes (MCC\,=\,$-$0.225), and handcrafted physicochemical terminal features (MCC\,=\,$-$0.065), all of which failed. Exploratory analyses using pre-specified protein-length windows suggested potentially elevated signal in moderate-length proteins (150--250\,aa, MCC\,=\,0.206) relative to shorter or longer proteins (MCC range $-$0.098 to 0.029), though this did not reach significance under permutation testing ($p$\,=\,0.118, $n$\,=\,44) and is reported as hypothesis-generating. Repeated discovery/validation splitting (10 independent 50:50 partitions) showed that decision-margin confidence was associated with higher prediction accuracy in 7 of 10 splits (mean accuracy 0.700 for high-confidence predictions versus 0.629 for low-confidence predictions), offering a practical, if modest, calibration heuristic. Only one of eight physicochemical features (disorder-prone fraction) showed a nominally significant difference between N- and C-terminal winners ($p$\,=\,0.037), which did not survive Bonferroni correction ($\alpha$\,=\,0.00625), confirming the benchmark presents a genuine predictive challenge rather than a trivially separable one. The TargetTrackHisTag dataset is deposited at Zenodo (DOI: [ZENODO\_DOI]).\\
\textbf{Availability:} Dataset and all analysis code: \url{https://github.com/uzairamu/HisTagPred}. Zenodo deposit: [ZENODO\_DOI].\\
\textbf{Contact:} \href{rizwanhkhan1@yahoo.com}{rizwanhkhan1@yahoo.com}}

\maketitle

\section{Introduction}

Recombinant protein production underpins modern biochemistry, structural biology, and biotechnology, with applications spanning enzyme engineering, drug discovery, vaccine development, and industrial biocatalysis \citep{burnett2020,rosano2019,overton2014}. Central to any recombinant expression workflow is purification of the target protein from complex cellular mixtures, for which affinity chromatography remains the method of choice owing to its high specificity, scalability, and compatibility with diverse expression hosts \citep{schmitt1993,li2010,bornhorst2000,zhao2013}.

The polyhistidine tag (His-tag) is the most widely employed affinity tag in recombinant protein production. Comprising a stretch of six to ten histidine residues, the His-tag forms stable complexes with divalent metal ions immobilised on chromatographic resin, enabling capture of His-tagged proteins from crude lysates with purities exceeding 80--95\% in a single step \citep{piatibratov2000,loughran2016}. The His-tag can be incorporated at either the N- or C-terminus of the target protein, and the choice of placement can have substantial consequences for expression yield, solubility, stability, and biological activity \citep{woestenenk2004,esposito2006,carson2007,kapust1999}.

Despite three decades of empirical evidence establishing that His-tag placement effects are protein-specific and unpredictable from sequence inspection alone \citep{waugh2011,terpe2003}, the determination of optimal placement remains dependent on parallel expression screening of both tagged variants --- a resource-intensive process that represents a recurring bottleneck in structural genomics pipelines and high-throughput protein production. The N-terminus and C-terminus of a protein are structurally and functionally distinct \citep{jacob2007}: the N-terminus participates in co-translational folding, may encode signal peptides or targeting sequences, and frequently contributes to early secondary structure formation \citep{netzer1997,fersht1999}, while the C-terminus commonly contributes to domain closure, intermolecular interactions, and catalytic function \citep{nooren2003,branden2012}. A His-tag inserted at a structurally constrained terminus --- one involved in folding nucleation, domain packing, or functional contacts --- is more likely to perturb protein behaviour than one placed at a disordered or solvent-exposed terminus \citep{wright1999,booth2018}.

With the emergence of protein language models (PLMs) and deep learning-based structure prediction, new possibilities have arisen for sequence-based prediction of protein properties that were previously inaccessible \citep{rives2021,lin2023}. ESM-2, a transformer-based PLM trained on 250 million sequences from UniRef, produces rich per-residue representations encoding evolutionary, structural, and physicochemical context, and has demonstrated strong performance across diverse protein property prediction tasks \citep{lin2023}. ESMFold further enables rapid single-sequence structure prediction, producing per-residue confidence scores (pLDDT) that correlate with local structural disorder and flexibility \citep{jumper2021}.

Despite these advances, two fundamental obstacles have prevented progress in computational His-tag placement prediction. First, no experimentally validated benchmark dataset exists against which computational methods can be rigorously evaluated --- prior work on recombinant protein expression optimisation has focused on solubility prediction, aggregation propensity, and expression level, without addressing the terminal-specific question of tag placement tolerance \citep{hebditch2017,sormanni2015,kuriata2019}. Second, the only structure-based method directly addressing this problem, Terminator \citep{gerulskis2025}, requires crystallographic structures and ConSurf conservation scores, limiting its applicability to the small fraction of proteins with solved structures.

We initially approached this problem by developing a sequence-based predictor using ESMFold-derived pseudo-labels and ESM-2 terminal embeddings. Under internal cross-validation, this approach appeared promising (XGBoost MCC\,=\,0.187, ROC-AUC\,=\,0.629). However, pseudo-labels derived from ESMFold pLDDT drop are themselves model-derived quantities rather than direct experimental measurements, and their correspondence to true experimental His-tag tolerance had not been established. To assess this correspondence, we sought an independent experimental benchmark against which the pseudo-label-based predictions could be evaluated. This validation revealed that the pseudo-label approach performed below random on experimental data (MCC\,=\,$-$0.243), indicating that internal cross-validation performance on pseudo-labels does not guarantee correspondence with biological reality, and motivating construction of a proper experimental benchmark.

Here we present TargetTrackHisTag, the first machine-readable benchmark of experimentally derived bilateral His-tag placement outcomes, curated from the NESG TargetTrack structural genomics archive. We describe the data mining pipeline, validation methodology, and systematic evaluation of sequence-based approaches against the benchmark. In contrast to our initial pseudo-label-based approach, full-sequence ESM-2 embeddings evaluated directly against this experimental benchmark achieve a statistically significant, if modest, predictive signal. We further characterise this signal through exploratory length-stratified analyses and confidence calibration, and discuss its implications alongside the limitations that motivate continued development of structure-aware approaches.

\section{Materials and Methods}

\subsection{TargetTrack Data Mining Pipeline}
The TargetTrack database \citep{xiao2010} is a publicly accessible repository maintained by the NESG that tracks experimental progress for protein targets through the structural genomics pipeline, from gene selection through cloning, expression, purification, and crystallisation. The complete TargetTrack XML archive (5.4\,GB) was downloaded and parsed using custom Python scripts. Entries were filtered to retain only His-tag-bearing constructs, yielding 97,418 relevant entries. For each target identifier, all recorded construct entries were retrieved, including sequence, construct orientation (N-terminal or C-terminal His-tag), and pipeline progression stage.

Progression stage was converted to a numerical score reflecting experimental advancement: selected (1), cloned (3), expressed (4), purified (5), crystals obtained (12), structure solved (15). For each target, the maximum score achieved across all N-terminal constructs was designated the N-score, and the maximum score across all C-terminal constructs was designated the C-score. Targets were assigned a winner label (N or C) based on which orientation achieved the higher maximum score; targets with equal scores were classified as ties and excluded. This scoring procedure yielded 804 candidate bilateral targets with non-tie outcomes.

\subsection{Alignment-Based Validation}
Initial sequence identity analysis of candidate pairs using naive position-by-position comparison yielded apparently low identities (3--5\%), incorrectly suggesting the N- and C-terminal constructs were unrelated proteins. Manual inspection revealed that these low identities arose from comparing tag-prefixed and tag-suffixed sequences without alignment: N-terminal constructs carry the His-tag as a prefix (e.g., MGHHHHHHSHM + protein) while C-terminal constructs carry it as a suffix (protein + LEHHHHHH), making direct comparison meaningless.

Following tag-stripping by removing canonical His-tag patterns from sequence termini using regular expression matching, pairwise alignment was performed for all 804 candidate pairs. Coverage was defined as the fraction of the shorter sequence aligned with the longer, after tag removal. Filtering at coverage $\geq$\,0.85 retained 156 proteins; filtering at coverage $\geq$\,0.95 retained 101 proteins. Relaxing the coverage threshold from 0.95 to 0.85 increased dataset size by 54.5\% (101 $\rightarrow$ 156 proteins) while preserving identical median score difference (delta\,=\,4.0) and only 4.3\% reduction in median confidence score, indicating the additional proteins represent genuine bilateral pairs rather than noise.

\subsection{Dataset Quality Audits}
Three automated audits were performed to assess label quality. Audit A verified tag pattern authenticity: for each protein, the full construct sequence was examined for canonical His-tag patterns at the expected terminus, including standard pET-vector His10-thrombin tags (e.g., MGHHHHHHSHM--- and MGSSHHHHHHSSGLVPRGSH--- prefixes) identified via terminal motif matching. Of 156 proteins, 99.4\% (155/156) showed genuine His-tag patterns; a single protein showed no detectable His-tag motif and was retained as the sole unclear case. Audit B assessed label robustness using tag length (6--25 residues, expected for His-tag plus linker) and score delta distributions; all 156 proteins (100\%) passed medium-to-high confidence thresholds. Audit C performed biological characterisation: Mann--Whitney U tests for eight physicochemical properties (sequence length, molecular weight, hydrophobicity, net charge, isoelectric point, N-terminal charge, C-terminal charge, disorder-prone fraction \citep{uversky2019}) between N-winners and C-winners. One feature (disorder-prone fraction) showed a nominally significant difference ($p$\,=\,0.037), which did not survive Bonferroni correction for multiple comparisons ($\alpha$\,=\,0.05/8\,=\,0.00625); the remaining seven features showed no significant differences (all $p$\,$>$\,0.16).

\subsection{ESM-2 Embedding Extraction}
Full-sequence ESM-2 embeddings were extracted for all 156 proteins using the esm2\_t33\_650M\_UR50D model (650 million parameters, 33 transformer layers, 1280-dimensional representations) from the fair-esm library with PyTorch. For each protein, the tag-stripped sequence was passed through the frozen model and per-residue representations from layer 33 were mean-pooled across all sequence positions to produce a single 1,280-dimensional embedding vector per protein.

\subsection{Machine Learning Evaluation Framework}
The primary evaluation used leave-one-out cross-validation (LOOCV) across the complete 156-protein dataset (108 C-winners, 48 N-winners), with logistic regression (C\,=\,0.1, class\_weight\,=\,balanced, max\_iter\,=\,2000) trained on standardised features (StandardScaler fit on training folds only) to account for class imbalance without discarding data through undersampling. Matthews Correlation Coefficient (MCC) was used as the primary performance metric owing to its robustness under class imbalance \citep{chicco2020}. This full-dataset evaluation was adopted as the primary analysis in preference to a balanced-undersampling approach (in which the majority class is randomly subsampled to match the minority class), which we also conducted for comparison; the two approaches yielded similar but not identical results (full dataset: MCC\,=\,0.168, $n$\,=\,156; undersampled: MCC\,=\,0.188, $n$\,=\,96), and we report the full-dataset result as primary since it does not discard experimental data and achieved statistical significance under permutation testing.

Statistical significance was assessed by permutation testing \citep{ojala2010}: 1,000 permutations of class labels were used to generate a null distribution of MCC values under the same LOOCV procedure, with the empirical $p$-value computed as the fraction of permuted MCCs equal to or exceeding the observed value.

\subsection{Length-Stratified Exploratory Analysis}
To assess whether predictive signal varies with protein length, we conducted four complementary analyses, each successively more conservative with respect to multiple-comparison concerns. First, proteins were divided into length quartiles and LOOCV MCC was computed within each quartile. Second, a series of cumulative length thresholds ($>$150, $>$175, $>$194, $>$225, $>$250, $>$275\,aa) was tested, with LOOCV MCC computed for proteins exceeding each threshold. Third, to assess whether any length-threshold effect would generalise to unseen data, the dataset was repeatedly partitioned into discovery and validation halves (10 independent 50:50 stratified splits); within each discovery half, the threshold from a pre-specified candidate set (150--225\,aa) achieving the highest LOOCV MCC was selected, and that exact threshold was then applied without further optimisation to the independent validation half. Fourth, to eliminate optimisation bias entirely, four fixed, pre-specified, non-overlapping length windows (50--150, 150--250, 250--350, and $>$350\,aa) were defined prior to any analysis of their performance, and LOOCV MCC was computed within each window, with permutation testing applied to the window of greatest a priori biological interest (150--250\,aa, corresponding to typical single-domain globular protein size).

\subsection{Confidence Calibration Analysis}
To assess whether the magnitude of the logistic regression decision function (decision margin) is informative of prediction reliability, we conducted ten independent 50:50 discovery/validation splits. Within each split, a model was trained on the discovery half and applied to the independent validation half; validation-half predictions were stratified into high-confidence (top tercile of absolute decision margin within the validation half) and lower-confidence (remaining two-thirds) groups, and accuracy was compared between groups.

\subsection{Pseudo-Label Validation}
To assess whether ESMFold-derived pseudo-labels correspond to experimental reality, the original XGBoost \citep{chen2016} model trained on 1,567 pseudo-labelled proteins (terminal ESM-2 2,560-dimensional embeddings, stratified 5-fold CV) was evaluated against a 28-protein experimental benchmark curated from published literature \citep{gerulskis2025,woestenenk2004}. ESMFold pLDDT-based pseudo-labels were also directly evaluated against experimental outcomes by computing pLDDT drop values for experimental proteins and comparing assigned pseudo-labels against known experimental preferences.

\subsection{Coarse-Grained Molecular Dynamics Evaluation}
To test whether structural dynamics rather than sequence or static structure encode the relevant terminal signal \citep{thielges2011}, we implemented an automated coarse-grained molecular dynamics pipeline (ESMFold structure prediction $\rightarrow$ Martini 3 Go-model coarse-graining via martinize2 \citep{souza2021,kroon2023,yang2025} $\rightarrow$ GROMACS \citep{berendsen1995} simulation) and evaluated seven geometric and dynamic descriptors (terminal RMSF with and without tag-bead exclusion, radius of gyration, terminal-to-core distance, terminal burial, terminal solvent exposure, and Go-contact stability) across a pilot set of 10 proteins with 5 independent replicas (100\,ns production simulation per replica) each. Terminal solvent exposure, the only descriptor showing replica-reproducible directional accuracy above chance in the pilot (6--7/10), was subsequently evaluated at scale using static ESMFold-structure-derived solvent-accessible surface area (SASA) asymmetry across the complete 156-protein benchmark and, independently, across the Gerulskis and Minteer Terminator benchmark \citep{gerulskis2025} ($n$\,=\,12 proteins with unambiguous experimental labels), to determine whether the metric generalised beyond the pilot set.

\subsection{Implementation}
All analyses were performed in Python 3.10. ESM-2 embeddings were extracted using fair-esm with PyTorch. Machine learning models were implemented using scikit-learn \citep{pedregosa2011} and XGBoost \citep{chen2016}. TargetTrack XML parsing used Python's xml.etree.ElementTree. Sequence alignment and tag-stripping used Biopython \citep{cock2009} and custom regular expressions. Coarse-grained simulations used GROMACS \citep{berendsen1995} with the Martini 3 force field and martinize2/vermouth for topology generation. Statistical tests used SciPy. All code is available at \url{https://github.com/uzairamu/HisTagPred}.

\section{Results}

\subsection{TargetTrackHisTag: Dataset Construction and Validation}
Parsing the complete TargetTrack XML archive (5.4\,GB, 97,418 His-tag entries) yielded 804 proteins with construct entries at both N- and C-terminal orientations. Following correction of the sequence identity analysis pipeline, 156 proteins passed the alignment coverage filter at $\geq$\,85\%. The final TargetTrackHisTag benchmark comprises 156 proteins (108 C-winners, 48 N-winners) with median alignment coverage 0.961, median score delta 4.0, and sequence lengths ranging from 59 to 622 amino acids (mean 218.3\,$\pm$\,115.3\,aa). Dataset characteristics are summarised in Table~\ref{tab:dataset}.

\begin{table}[!t]
\processtable{TargetTrackHisTag benchmark characteristics\label{tab:dataset}}
{\begin{tabular}{ll}
\toprule
Parameter & Value\\
\midrule
Total proteins (coverage $\geq$ 0.85) & 156\\
C-terminal winners & 108 (69.2\%)\\
N-terminal winners & 48 (30.8\%)\\
Mean sequence length & 218.3 $\pm$ 115.3\,aa\\
Sequence length range & 59--622\,aa\\
Median alignment coverage & 0.961\\
Median score delta & 4.0\\
Tag pattern verification (Audit A) & 99.4\% (155/156)\\
Label confidence (Audit B) & 156/156 (100\%)\\
Sig.\ N vs C features (Audit C) & 1/8 nominal; 0/8 Bonf.\\
Zenodo DOI & [ZENODO\_DOI]\\
\botrule
\end{tabular}}{}
\end{table}

\subsection{Pseudo-Label Validation: ESMFold pLDDT Drops Do Not Correspond to Experimental Reality}
The initial prediction framework, trained on ESMFold pLDDT-derived pseudo-labels for 1,567 proteins, achieved promising internal cross-validation performance (XGBoost MCC\,=\,0.187\,$\pm$\,0.038, ROC-AUC\,=\,0.629\,$\pm$\,0.016). However, evaluation against the 28-protein experimental benchmark revealed near-total failure: the pseudo-label XGBoost model achieved MCC\,=\,$-$0.225 on experimental data, performing below random. Direct evaluation of ESMFold pLDDT drops against experimental outcomes also yielded below-random performance (MCC\,=\,$-$0.243, accuracy\,=\,40.0\%). These results indicate that pLDDT-based confidence metrics do not encode the biophysical signal determining experimental His-tag tolerance.

\subsection{Full-Sequence ESM-2 Embeddings Achieve a Significant Signal on the Experimental Benchmark}
In contrast to the pseudo-label approach and to handcrafted physicochemical terminal features (MCC\,=\,$-$0.065, $n$\,=\,62 balanced subset), full-sequence ESM-2 embeddings evaluated by logistic regression across the complete 156-protein benchmark achieved MCC\,=\,0.168 (balanced accuracy\,=\,0.583; confusion matrix: 81 true C, 27 false N, 28 false C, 20 true N). Permutation testing (1,000 permutations) confirmed this result exceeds chance expectation (observed MCC\,=\,0.168, null mean\,=\,$-$0.001, null SD\,=\,0.105, $Z$\,=\,1.60, $p$\,=\,0.0498). An alternative evaluation using balanced undersampling ($n$\,=\,96, equal C and N representation) yielded a similar MCC (0.188, balanced accuracy\,=\,0.594) but did not reach significance under the same permutation procedure ($p$\,=\,0.084), illustrating that retaining the full, imbalanced dataset under appropriate class weighting provided greater statistical power than artificial balancing through undersampling. We report the full-dataset result (MCC\,=\,0.168, $p$\,=\,0.0498) as the primary finding. This is, to our knowledge, the first statistically significant sequence-based predictive signal for His-tag placement preference reported against an experimentally derived benchmark. Figure~\ref{fig:confusion} shows the corresponding confusion matrix.

\begin{figure}[!t]
\centerline{\includegraphics[width=0.45\textwidth]{Figure1_confusion_matrix_v3.png}}
\caption{Confusion matrix for the full-dataset ESM-2 logistic regression model ($n$=156, LOOCV, MCC=0.168). True C=81, False N=27, False C=28, True N=20, demonstrating bidirectional, non-majority-collapsed prediction.}\label{fig:confusion}
\end{figure}

\subsection{Exploratory Length-Stratified Analysis: A Suggestive but Non-Significant Pattern}
Given the modest effect size of the full-dataset result, we explored whether predictive signal varies systematically with protein length. Quartile-based analysis (Q1: 59--134\,aa, MCC\,=\,0.112; Q2: 139--194\,aa, MCC\,=\,$-$0.218; Q3: 195--292\,aa, MCC\,=\,0.355; Q4: 294--622\,aa, MCC\,=\,0.120) did not reveal a monotonic trend. A median-length split ($>$194\,aa, $n$\,=\,78) yielded a higher apparent signal (MCC\,=\,0.339, balanced accuracy\,=\,0.666) than the unrestricted dataset, which itself reached significance under a subset-specific permutation test ($Z$\,=\,2.62, $p$\,=\,0.0080); however, a Spearman correlation between protein length and the magnitude of the model's decision margin across the full dataset was weak and in the direction opposite to the length-dependence hypothesis ($r$\,=\,$-$0.170, $p$\,=\,0.0343), and a sliding-window analysis (window\,=\,40 proteins) revealed substantial non-monotonic fluctuation rather than a smooth trend (MCC ranging from $-$0.246 to 0.217 across overlapping windows).

To assess whether an optimal length threshold would generalise to independent data, we repeated a discovery/validation procedure across ten independent 50:50 stratified splits, selecting the best-performing threshold within each discovery half and applying it without further optimisation to the corresponding validation half. Validation-half MCC was positive in 8 of 10 splits (mean\,=\,0.166\,$\pm$\,0.181), comparable in magnitude to the full-dataset MCC (0.168) rather than exceeding it, indicating that length-based restriction redistributes rather than amplifies the overall signal. Finally, four fixed, non-overlapping length windows defined prior to outcome inspection (50--150, 150--250, 250--350, and $>$350\,aa) yielded MCC values of 0.014, 0.206, 0.029, and $-$0.098 respectively; permutation testing of the window with the highest observed MCC (150--250\,aa, $n$\,=\,44) did not reach significance ($Z$\,=\,1.26, $p$\,=\,0.118). We therefore report the length-dependence pattern as a suggestive but statistically unconfirmed observation, consistent across several complementary analyses in direction but not robust to permutation testing of a pre-specified window. Figure~\ref{fig:length} summarises the length-stratified and confidence-calibration analyses.

\begin{figure}[!t]
\centerline{\includegraphics[width=0.49\textwidth]{Figure5_length_confidence.png}}
\caption{Length-stratified and confidence-calibration analysis. (A) MCC across quartile length bins. (B) Discovery/validation MCC across 10 independent splits. (C) Confidence-stratified accuracy (high vs.\ lower decision-margin tercile) across the same 10 splits.}\label{fig:length}
\end{figure}

\subsection{Confidence-Stratified Prediction Accuracy}
Across the same ten discovery/validation splits used for the length analysis, predictions on the validation half were stratified by decision-margin confidence. High-confidence predictions (top tercile of absolute decision margin) achieved higher mean accuracy than lower-confidence predictions in 7 of 10 splits (overall mean accuracy: 0.700 for high-confidence versus 0.629 for lower-confidence predictions). While this difference did not reach formal statistical significance given the limited number of independent splits, the consistent direction and magnitude of the effect across most partitions suggest that decision-margin magnitude may offer a practical, if modest, heuristic for prioritising more reliable predictions in applied use.

\subsection{Comparison with Structure-Based Prediction and Coarse-Grained Molecular Dynamics}
For context, Terminator \citep{gerulskis2025} achieved 86--93\% accuracy on 14 proteins using cavity proximity scores derived from crystallographic structures combined with ConSurf \citep{ashkenazy2016} evolutionary conservation scores. To test whether equilibrium structural dynamics, rather than sequence, might capture complementary or superior signal, we implemented an automated coarse-grained molecular dynamics pipeline and evaluated seven terminal geometric and dynamic descriptors across a 10-protein pilot with five independent replicas per protein. No descriptor achieved replica-robust directional accuracy clearly above chance; the most promising candidate, terminal solvent exposure (directionally consistent across replicas in 7 of 10 pilot proteins), was subsequently evaluated at scale using static ESMFold-derived solvent-accessible surface area asymmetry across the complete 156-protein benchmark (AUC\,=\,0.456, 95\% CI [0.355, 0.558], below the 69.2\% majority-class baseline) and independently against the Terminator benchmark (AUC\,=\,0.074, $n$\,=\,12), confirming the pilot-stage result did not generalise. These results indicate that neither equilibrium terminal dynamics nor static structural accessibility, whether derived from coarse-grained molecular dynamics or ESMFold structures, captures additional signal beyond what is already present in ESM-2 sequence embeddings, reinforcing that the moderate sequence-based signal we report and the cavity-proximity signal exploited by Terminator likely reflect distinct, only partially overlapping, biophysical determinants. Figure~\ref{fig:comparison} summarises the full comparative evaluation, and Figure~\ref{fig:audits} summarises the dataset quality audits.

\begin{figure}[!t]
\centerline{\includegraphics[width=0.49\textwidth]{Figure3_comparison_v3.png}}
\caption{Systematic evaluation against experimental His-tag placement benchmarks. MCC across all evaluated approaches, including the significant full-dataset ESM-2 result, pseudo-label and structural-confidence failures, exploratory length windows, and the coarse-grained MD/static-SASA evaluation.}\label{fig:comparison}
\end{figure}

\begin{figure}[!t]
\centerline{\includegraphics[width=0.49\textwidth]{Figure4_audits.png}}
\caption{Dataset quality audit results. (A) Tag pattern classification (Audit A). (B) Label confidence distribution (Audit B). (C) Physicochemical feature comparison (Audit C), with Bonferroni threshold indicated.}\label{fig:audits}
\end{figure}

\section{Discussion}

This study presents TargetTrackHisTag, the first machine-readable experimental benchmark for His-tag placement prediction, and reports the first statistically significant sequence-based predictive signal for this problem evaluated against experimentally derived labels (MCC\,=\,0.168, permutation $p$\,=\,0.0498, $n$\,=\,156). This result stands in contrast to our initial pseudo-label-based approach, which appeared promising under internal cross-validation but failed when evaluated against experimental data (MCC\,=\,$-$0.225), and to handcrafted physicochemical features and structural confidence metrics, all of which showed no signal. The magnitude of the validated signal is modest --- comparable to weak-to-modest effect sizes reported for other challenging sequence-based prediction tasks --- and we are careful throughout to avoid overstating its practical utility as a standalone deployable predictor.

The failure of pseudo-label approaches provides an important methodological lesson that extends beyond this specific problem. ESMFold pLDDT drop appeared to be a biologically motivated labelling criterion, and internal cross-validation performance (MCC\,=\,0.187) appeared to support this. However, validation against experimental data revealed that this signal was circular: both the labels and the predictive features derived from related model outputs, and both failed when confronted with independent experimental outcomes. The general lesson is that internal cross-validation performance on pseudo-labels is not a reliable indicator of biological validity \citep{natarajan2013}, and experimental benchmarks are essential for evaluating computational protein property predictors, particularly those relying on pseudo-labelling strategies increasingly common in computational biology.

We undertook an extensive, multi-pronged exploration of whether the modest overall signal might be concentrated within identifiable protein subpopulations, motivated by the observation that an initial median-length split appeared to show a substantially stronger effect (MCC\,=\,0.339) than the unrestricted dataset. Systematic testing across quartile analysis, cumulative threshold scanning, repeated discovery/validation splitting, and pre-specified fixed-window analysis consistently showed a directionally favourable but statistically unconfirmed pattern: proteins of moderate length (approximately 150--250\,aa) tended to show numerically higher MCC than very short or very long proteins, but no single length threshold proved stable across independent validation splits, a Spearman correlation between length and prediction confidence trended in the opposite direction, and permutation testing of the most promising pre-specified window did not reach significance ($p$\,=\,0.118). We report this length-dependence as hypothesis-generating rather than confirmed, and explicitly caution against its use as the basis for a length-restricted deployable predictor without independent experimental validation. We consider this extended negative-but-thorough exploration scientifically valuable in its own right: it demonstrates that an initially striking subgroup finding, of the kind frequently reported in smaller or less rigorously validated studies, did not survive appropriately conservative statistical scrutiny, and we present the complete analytical trail to allow readers to evaluate this judgement directly.

The confidence-calibration analysis offers a more practically actionable, if still modest, finding: predictions made with larger decision-margin magnitude were more often correct than lower-margin predictions across most (7 of 10) independent validation splits. This is consistent with standard expectations for probabilistic classifiers and suggests that, even though the overall predictor is not accurate enough for unconditional deployment, a practical workflow restricting experimental follow-up to high-confidence predictions may offer modest efficiency gains over uninformed selection between N- and C-terminal constructs.

The comparison with Terminator (86--93\% accuracy using crystal structures) against the considerably more modest performance of sequence-based approaches reported here suggests that the biophysical signal most relevant to experimental success --- plausibly related to steric obstruction of terminal cavities by the His-tag --- is only partially accessible from sequence alone. Our systematic evaluation of coarse-grained molecular dynamics and static structural accessibility metrics, both of which failed to generalise beyond an encouraging pilot-stage result, further suggests that equilibrium terminal dynamics and static solvent accessibility, whether derived from physics-based simulation or from predicted structure, do not capture the relevant signal either. Cavity proximity within quaternary assemblies, as exploited by Terminator using crystallographic data, may represent a determinant not accessible to any of the monomer-level computational approaches --- sequence-based, dynamics-based, or static-structure-based --- evaluated in this study.

Several limitations of TargetTrackHisTag should be acknowledged. The 70:30 class imbalance (C-winners:N-winners) reflects empirical preference for C-terminal His-tags in the NESG pipeline and may not represent the full distribution of His-tag tolerance in the broader proteome. The labelling scheme assigns winners based on maximum pipeline progression score, which conflates experimental success with factors unrelated to tag placement. The 156-protein dataset, while the largest experimental His-tag placement benchmark assembled to date, remains modest relative to the dimensionality of sequence-based features, limiting statistical power, particularly for subgroup analyses as demonstrated by our length-stratification results. The NESG dataset is dominated by bacterial and archaeal proteins selected for structural genomics utility, which may not be representative of eukaryotic protein families most relevant to pharmaceutical and biotechnology applications.

Future development of TargetTrackHisTag should prioritise dataset expansion through mining additional structural genomics archives (JCSG, MCSG, SGC) and recent high-throughput expression screening studies, which would directly address the statistical power limitations encountered in our length-stratified and confidence-calibration analyses and would allow more confident resolution of whether the length-dependence pattern we observed reflects a genuine, if modest, biological phenomenon or residual noise. Development of structure-aware methods using AlphaFold-Multimer predictions \citep{evans2022} combined with evolutionary conservation proxies, or with structure-based stability and binding-energy estimation tools such as FoldX \citep{schymkowitz2005}, could extend Terminator-like approaches beyond proteins with known crystal structures, addressing the most practically important limitation of current structure-based methods.

In summary, TargetTrackHisTag provides the computational biology community with the first experimentally grounded benchmark for His-tag placement prediction and establishes, for the first time, a statistically significant though modest sequence-based predictive signal. We additionally demonstrate, through extensive and statistically conservative exploration, that this signal does not appear to be concentrated within a robustly identifiable length-restricted subpopulation, and that decision-margin confidence offers a practical heuristic for prioritising more reliable predictions. These results motivate continued benchmark expansion and development of structure-aware methods capable of capturing the cavity-proximity signal that current sequence-based and dynamics-based approaches do not access.

\section*{Acknowledgements}
The authors thank the Interdisciplinary Biotechnology Unit, AMU, for providing computational resources.

\section*{Funding}
This study did not receive external funding. M.U.A.\ is supported by the Prime Minister's Research Fellowship (PMRF ID: 3802926).

\section*{Conflict of interest}
None declared.

\section*{Author contributions}
M.U.A.: Conceptualisation, Data curation, Formal analysis, Investigation, Methodology, Project administration, Software, Validation, Writing -- original draft, Writing -- review and editing. Muazzam: Data curation, Methodology, Resources, Software. R.H.K.: Project administration, Supervision, Validation, Writing -- review and editing.

\section*{AI tool usage declaration}
During the preparation of this work, the authors used Claude (Anthropic) to assist with manuscript drafting, code development, and language editing. After using this tool, the authors reviewed and edited all content as needed and take full responsibility for the content of the publication.

\section*{Data availability}
The TargetTrackHisTag benchmark dataset (experimental\_156\_clean.csv, target\_ids\_156.csv, audit results, length-stratified and confidence-calibration analysis outputs) is deposited at Zenodo (DOI: [ZENODO\_DOI]) under CC BY 4.0 licence. All analysis code, the coarse-grained molecular dynamics pipeline, ESM-2 embedding extraction scripts, and LOOCV evaluation pipelines are available at \url{https://github.com/uzairamu/HisTagPred}.

\bibliographystyle{oup-abbrvnat}
\bibliography{references}

\end{document}
